% AI-produced by GPT-6 Astra (OpenAI, ChatGPT Pro mode). No priority claim.
\subsection{A source-specific local convexity bound}

The counting invariant of C.2 gives a uniform local geometric statement
which is stronger than inducedness and uses no extra choice of target
size.  It also identifies a class of source objects for which the
fixed-path kernel vanishes at every later stage.

\begin{proposition}[Proved: the shortest possible outside excursion]
\label{C:local-convexity}
Let $X\in K_m$, let $h:Y\to X$ be an arrow of the repaired C.1
construction, and let $J=h_\#G_Y\subseteq G_X$.  If a simple graph
path has distinct endpoints in $J$, all its internal vertices outside
$J$, and all its edges outside $J$, its length $r$ satisfies
\[
             (2-a_m)r\geq1.
\]
For $m\geq1$ put
\[
       R_m=\left\lceil\frac1{2-a_m}\right\rceil-1.
\]
Every simple path of length at most $R_m$ with endpoints in $J$
is contained in $J$.  In particular every arrow image is
\emph{two-convex}: every geodesic of length at most two whose
endpoints lie in that image stays in the image.  This holds at all
stages and for every admissible numerical choice.
\end{proposition}

\begin{proof}
The graph $G_X$ is simple by $(c_m)$ and
$\operatorname{End}(W)=\{1\}$.  Its vertices are the underlying
points of $X$, as in Proposition~\ref{C:fixed-path-kernel}.
Write $b_1,\ldots,b_r:W\to X$ for the distinct arrows supplying
the edges of the displayed simple excursion.  None factors through
$h$, because its binary edge is outside $J$.  Invariant $(c_m)$
therefore gives $|b_i(W)\cap h(Y)|\leq1$ for every $i$.
The relative counting invariant $(d_m)$ applies and yields
\[
       \left|\bigcup_{i=1}^r b_i(W)\setminus h(Y)\right|
                      \geq a_m r.
\]
The binary endpoints of these edges contribute exactly the $r-1$
internal path vertices outside $h(Y)$.  Each $b_i$ has only one
further point, namely $b_iv_3$, and these contribute at most $r$
additional outside points.  Consequently the same union has size
at most $2r-1$.  Combining the bounds proves
$a_m r\leq2r-1$.

If a simple path with endpoints in $J$ leaves $J$, take a maximal
consecutive excursion between two successive visits to $J$.
Its endpoints are distinct, because the entire path is simple.
Its edges are outside $J$ and its internal vertices are outside
$J$.  Thus its length is at least $1/(2-a_m)$, and the length of
the original path is at least that number as well.  This proves
the bound with $R_m$.  Since $a_m>3/2$, one has
$1/(2-a_m)>2$, hence $R_m\geq2$.  At stage zero the assertion is
immediate from the singleton, edge, and isolated-vertex graphs;
equivalently $a_0=2$ makes the displayed excursion inequality
impossible.
\end{proof}

\begin{corollary}[Proved: small-diameter sources never kill a class]
\label{C:small-diameter-source}
Let $f,g:S\rightrightarrows T$ be parallel arrows at an old stage
$K_m$, where $m\geq1$.  If each component of $G_S$ has diameter at
most $R_m$, then its equalizer subgraph $F$ is connected in every
component of $G_S$ that it meets.  A repaired quasi-filter attachment
for this pair therefore preserves the whole old colimit of $K$,
and in particular the whole old binary second homology.

The same conclusion holds at any later stage whenever every
component of the fixed source $G_S$ has diameter at most two.
It does not require prior component injectivity of $f$ or $g$.
\end{corollary}

\begin{proof}
Let $u,v$ be fixed vertices in one source component and put
$d=d_{G_S}(u,v)\leq R_m$.  Each embedding carries a source shortest
path to a target path of length $d$.  A target shortest path between
the image endpoints therefore has length at most $d$, and
Proposition~\ref{C:local-convexity} puts every such path in each
arrow image.  Inducedness and injectivity then lift it to a source
path.  In particular the target distance equals $d$.

Let $I_S(u,v)$ be the finite set of all source vertices on shortest
$u$--$v$ paths, and write $I_T$ for the corresponding target interval
between the common images of $u,v$.  The preceding argument gives
\[
                         f(I_S(u,v))=I_T=g(I_S(u,v)).
\]
Hence $g^{-1}f$ is a permutation of $I_S(u,v)$.  A nonidentity
cycle of this permutation would give a nontrivial directed cycle
of the point-transition relation $f(x)\mapsto g(x)$, contrary to
the established AKR bounded-chain property (4).  This permutation
is therefore the identity.  Every vertex of every shortest source
path from $u$ to $v$ is fixed.  The equalizer graph is induced on
its fixed vertices by Proposition~\ref{C:fixed-path-kernel}, so
these paths lie in $F$.

Thus $\pi_0(F)\to\pi_0(G_S)$ is injective.  The exact fixed-path
kernel formula proves the preservation assertion.  Finally
$R_m\geq2$ at every stage, giving the uniform diameter-two case.
\end{proof}

\begin{theorem}[Proved: every first-stage source remains harmless]
\label{C:first-stage-safe-sources}
In any repaired C.1 construction, let $S$ be an object already in
$K_1$.  At any later finite stage, every repaired quasi-filter
attachment for any pair $f,g:S\rightrightarrows T$ preserves the
whole old colimit of the binary cycle functor $K$.  It therefore
preserves every old binary second-homology class, independently of
the target $T$, the chosen pair, and all admissible numerical
parameters.

Consequently an attachment with a nonzero fixed-path kernel, if one
occurs, must have a source object born at stage two or later.
\end{theorem}

\begin{proof}
We first classify the binary graphs of the first-stage objects,
allowing different admissible necklace lengths at their births.
The original objects, the targets for identity pairs, and all
first-stage disjoint sums have graph components of diameter at most
two.  More explicitly, the equal pair $(v_1,v_1)$ or $(v_2,v_2)$
produces a star together with isolated third points; the equal pair
$(v_3,v_3)$ produces disjoint edges and one isolated fixed point.
The identity pairs on $V$ or $W$ merely reproduce their underlying
graphs on a fresh target.

For a distinct pair involving $v_3$ and one binary endpoint, every
binary edge has one active endpoint on a repaired two-vertex point
path and one private endpoint.  The formulas for $\beta_i^2(0)$
and $\beta_i^2(1)$ show that each active point is used as that
binary endpoint at most twice: one coordinate is used twice, one
is omitted, and all remaining coordinates are used once.
Accordingly all nontrivial binary components are edges or
two-edge stars.  They have diameter at most two.  The only remaining
objects are the two protected targets $P_+,P_-$: each has its unique
circle of length $s_\pm\geq45$, with all other vertices isolated.

Corollary~\ref{C:small-diameter-source} handles every case except
$S=P_+$ or $S=P_-$.  In one of these two cases, if $f=g$ the entire
graph is fixed.  Suppose $f\ne g$.  Two fixed marked vertices
which were nonadjacent on the protected circle would lie in
$f(S_\pm)\cap g(S)$.  Proposition~\ref{C:protected-nonadjacent}
would factor $f$ through $g$.  The factor would be an endomorphism
of $S$, and the only such endomorphism is the identity.  This
would contradict $f\ne g$.

Thus the fixed subset of the circular component is empty, a
singleton, or a pair of adjacent vertices.  In the last case its
edge is fixed as well, since the equalizer graph is induced.
Every other source component is an isolated vertex.  The
equalizer graph is therefore connected in each source component
that it meets.  Proposition~\ref{C:fixed-path-kernel} again proves
that the attachment kernel is zero.
\end{proof}

\paragraph{Scope.}
The local bound and the first-stage safe-source theorem apply to
the arrows actually generated by C.1 and use the source's relative
counting invariant. They do not by themselves prove convexity for
arbitrary long paths, or connectedness of fixed sets for all later
source objects. The separate birth-cover and component-path
induction in Theorem~\ref{C:universal-components} supplies the
latter conclusion, and Corollary~\ref{C:universal-nontotal}
completes the universal-choice assertion in C.1.
